\section{来源审计与证据边界：先说明 GLOM 还没有“成功”}

这节课最容易被写坏的地方，不是公式，而是证据等级。Geoffrey Hinton 介绍的 \term{GLOM} 不是一个已经完成训练、在标准数据集上取得分数的模型；它首先是一份关于表示方式的设计文档：如果神经网络不能像程序一样动态申请一块内存来存 parse-tree node，那么能否用固定硬件中的动态 activity pattern 表示每张图像都不同的层次结构？本讲的价值在于把 capsule、contrastive learning、attention、distillation 与 neural field 放进一个统一假说，而不是给出已验证的产品配方。

\lecturefigure{slide-01-title.jpg}{讲座主题：如何在神经网络中表示 part-whole hierarchy。}{Stanford Online 官方视频 00:00:26--00:01:07。}

\begin{warningbox}{最重要的历史边界：这是 design sketch，不是 benchmark result}
Hinton 在开场直接把完整系统称为 imaginary system 和“vaporware”。原论文 \href{https://arxiv.org/abs/2102.12627}{arXiv:2102.12627} 共 44 页，公开记录只有 2021-02-25 的一个版本。课堂提到少量 toy fragments，但没有报告端到端 GLOM 在 segmentation、recognition 或 reasoning benchmark 上的完整结果。因此下文所有“应该形成 islands”“可以用 consensus 训练”等句子都表示设计机制与可检验假设，而不是已经证明的事实。
\end{warningbox}

\teachervoice{讲者说自己本来在写一个系统设计文档，最后发现设计文档本身已经足够有趣；“little bits”可能存在，但整套东西不存在。保留这句自我限制非常关键，因为它决定我们应该用 architecture review 的方式阅读，而不是用 leaderboard review 的方式阅读。}

\subsection{三类已有思想，压成一个表示问题}

GLOM 的路线不是发明一个全新的 layer，而是重新组合三个方向：Transformer 提供基于相似度的交互；unsupervised/contrastive representation learning 提供“哪些 activity 应当一致”的学习信号；neural field 提供“同一个高层对象如何在不同坐标位置生成不同低层部件”的位置条件函数。Capsule、distillation 和 recurrent settling 则贯穿其间。

\lecturefigure{slide-02-three-advances.jpg}{GLOM 试图组合 Transformer、对比表示学习与 neural field。}{Stanford Online 官方视频 00:01:08--00:01:51。}

\begin{knowledgebox}{本讲使用的证据类型}
\begin{tabularx}{\textwidth}{L{0.22\textwidth}L{0.31\textwidth}>{\raggedright\arraybackslash}X}
\toprule
证据类型 & 本讲能提供什么 & 不能直接推出什么 \\
\midrule
表示论动机 & 固定连接与动态 parse structure 的矛盾 & GLOM 一定可训练、一定优于现有模型 \\
心理演示 & 人类会选择 intrinsic coordinate frame 与 alternative parse & 大脑必然实现 GLOM 的具体向量更新式 \\
架构草图 & column、level、attention、top-down/bottom-up 的信息流 & 稳定收敛、样本效率、算力成本的实测数字 \\
与已有论文类比 & SimCLR、BERT、Hough、NeRF 提供可复用 mechanism & 这些系统的成功会自动转移到 GLOM \\
\bottomrule
\end{tabularx}
\end{knowledgebox}

\subsection{工程与脑科学不是同一个验收标准}

Hinton 把神经网络研究分成两类目标：工程目标只问系统能否工作，脑科学目标还问这种机制是否可能解释 biological intelligence。共享权重、百层 ResNet 对工程师完全合理，但未必是脑的直接模型。GLOM 故意同时触碰两类问题，因此讲义必须把 engineering plausibility 与 biological plausibility 分开。

\lecturefigure{slide-03-two-research-goals.jpg}{两种研究目标：做出有效技术，或借人工网络理解大脑。}{Stanford Online 官方视频 00:01:52--00:02:37。}

\teachervoice{讲者提醒：神经网络研究曾经长期依赖一个朴素信念继续存在——大脑已经证明复杂学习在原则上可能。这个历史判断不是说工程模型必须复制大脑，而是说 biological existence proof 仍可启发新的 representation question。}

\subsection{本章小结}

本讲讨论的是一个历史上明确受限的设计提案。阅读顺序应是：先理解它要解决的 representational contradiction，再检查每个机制是否形成可计算、可训练、可证伪的假说，最后才讨论它与现代 object-centric model 或 world model 的关系。

